% command.tex — macros for vldbNuclearR1
% Mirrors NuclearCD/command.tex (slim subset)

% --- inline-heading macros (NuclearCD style) ---
\newcommand{\stitle}[1]{{\noindent\bf{#1}}}
\newcommand{\sstitle}[1]{{\noindent\textit{\underline{#1}}}}
\newcommand{\ssstitle}[1]{{\noindent\textit{#1}}}

% --- algorithm name macros (use \kw for sans-serif keyword style) ---
\newcommand{\kw}[1]{{\ensuremath{\mathsf{#1}}}\xspace}
\newcommand{\bkw}[1]{{\ensuremath{\mathsf{\textbf{#1}}}}\xspace}

\newcommand{\cpi}{\kw{CPI}}
\newcommand{\sct}{\kw{SCT}}
% --- main core-value algorithms ---
% Umbrella family name SPIN (Static Path Index for Nucleus decomposition):
%   \spin     = direct fixed-point iteration over the static CPI
%   \spinStar = the same fixed point computed by single-pass bucket-queue peeling
\newcommand{\spin}{\kw{SPIN}}
\newcommand{\spinStar}{\ensuremath{\mathsf{SPIN}^{\star}}\xspace}
\newcommand{\paraBuild}{\kw{ParaBuild}}
\newcommand{\buildHier}{\kw{BuildHier}}
% Per-vertex / per-path support helpers (companion pair, same family).
\newcommand{\vertexSupport}{\kw{VertexSupport}}
\newcommand{\pathSupport}{\kw{PathSupport}}
% Comparison baseline (general (r,s)-nucleus decomposition restricted to r=1)
\newcommand{\cbnd}{\kw{CND}}

% --- path notation (CPI paths denoted by \TPath, takes optional subscript) ---
\makeatletter
\def\TPath{\@ifnextchar\bgroup{\TPathWithArg}{\TPathNoArg}}
\def\TPathWithArg#1{\mathcal{P}_{#1}}
\def\TPathNoArg{\mathcal{P}}
\makeatother

% V_h / V_p as standalone subscript symbols; written inline as V_h(\TPath)
\newcommand{\Vh}{V_h}
\newcommand{\Vp}{V_p}

% --- support / index size ---
\newcommand{\Sig}{\Sigma}
\newcommand{\sdeg}{\deg^{(s)}}     % s-clique support of a vertex
\newcommand{\kmax}{\kappa_{\max}}

% --- theorem-like environments (independent counters; reset per section) ---
\theoremstyle{plain}
\newtheorem{theorem}{Theorem}[section]
\theoremstyle{plain}
\newtheorem{lemma}{Lemma}[section]
\theoremstyle{plain}
\newtheorem{corollary}{Corollary}[section]
\theoremstyle{definition}
\newtheorem{definition}{Definition}[section]
\theoremstyle{remark}
\newtheorem{example}{Example}[section]
\theoremstyle{remark}
\newtheorem{observation}{Observation}[section]

% --- reference shorthand ---
\newcommand{\refsec}[1]{Section~\ref{sec:#1}}
\newcommand{\reffig}[1]{Fig.~\ref{fig:#1}}
\newcommand{\reftable}[1]{Table~\ref{tab:#1}}
\newcommand{\refalg}[1]{Algorithm~\ref{alg:#1}}
\newcommand{\refdef}[1]{Definition~\ref{def:#1}}
\newcommand{\refthm}[1]{Theorem~\ref{thm:#1}}
\newcommand{\reflem}[1]{Lemma~\ref{lem:#1}}
\newcommand{\refex}[1]{Example~\ref{ex:#1}}

% --- hold/pivot styling (sans-serif, not bold) ---
\newcommand{\hold}{\textsf{hold}\xspace}
\newcommand{\pivot}{\textsf{pivot}\xspace}
\newcommand{\holds}{\textsf{holds}\xspace}
\newcommand{\pivots}{\textsf{pivots}\xspace}

% --- algorithm pseudocode font size (matches NuclearCD) ---
\newcommand{\algfontsize}{\fontsize{7pt}{7.5pt}\selectfont}
